\subsection{环面的结点群}

环面 S 的结点群记作 \(\Gamma(\mathrm{S})\) ，它是指数映射 \(\mathrm{L}(\mathrm{S}) \to \mathrm{S}\) 的核。这是 \(\mathrm{L}(\mathrm{S})\) 的离散子群，其秩等于 S 的维数，且 R-线性映射 \(\mathbf{R} \otimes_{\mathbf{Z}} \Gamma(\mathrm{S}) \to \mathrm{L}(\mathrm{S})\) （它把 \(\Gamma(\mathrm{S})\) 到 \(\mathrm{L}(\mathrm{S})\) 中的典则单射加以扩张）是双射。它经过渡到商诱导一个同构 \(\mathbf{R}/\mathbf{Z} \otimes_{\mathbf{Z}} \Gamma(\mathrm{S}) \to \mathrm{S}\) 。

例如，U 的结点群 \(\Gamma(\mathbf{U})\) 是子群 \(2\pi iZ\) （\(\mathrm{L}(\mathbf{U}) = i\mathbf{R}\) 的）。

对环面的任一态射 \(f: \mathrm{S} \to \mathrm{S}'\) ，用 \(\Gamma(f)\) 表示同态 \(\Gamma(\mathrm{S}) \to \Gamma(\mathrm{S}')\) （由 \(\mathrm{L}(f)\) 诱导的）。我们有交换图式

\[
\begin{array}{c c c c c c c c c} 0 & \longrightarrow & \Gamma (\mathrm{S}) & \longrightarrow & \mathrm{L} (\mathrm{S}) & \stackrel {{\exp_ {\mathrm{S}}}} {{\longrightarrow}} & \mathrm{S} & \longrightarrow & 0 \\ & & \Big \downarrow^ {\Gamma (f)} & & \Big \downarrow^ {\mathrm{L} (f)} & & \Big \downarrow^ {f} \\ 0 & \longrightarrow & \Gamma (\mathrm{S} ^ {\prime}) & \longrightarrow & \mathrm{L} (\mathrm{S} ^ {\prime}) & \stackrel {{\exp_ {\mathrm{S} ^ {\prime}}}} {{\longrightarrow}} & \mathrm{S} ^ {\prime} & \longrightarrow & 0. \end{array}\tag{1}
\]

设 \(a \in \mathrm{X}(\mathrm{S})\) ；把上述应用到由 \(a\) 定义的从 S 到 U 的态射，我们看到第 1 节的 C-线性映射 \(\delta(a) : \mathrm{L}(\mathrm{S})_{(\mathbf{C})} \to \mathbf{C}\) 把 \(\Gamma(\mathrm{S})\) 映到 \(2\pi i\mathbf{Z}\) 中。于是，我们可以在 \(\mathrm{X}(\mathrm{S}) \times \Gamma(\mathrm{S})\) 上定义一个 Z-双线性型，方法是令

\[
\langle a, X \rangle = \frac {1}{2 \pi i} \delta (a) (X), \quad a \in \mathrm{X(S)}, X \in \Gamma (\mathrm{S}).\tag{2}
\]

命题 3. 双线性型 \((a,X)\mapsto \langle a,X\rangle\) （在 \(\mathrm{X(S)}\times \Gamma (\mathrm{S})\) 上）是可逆的。

回忆（《代数学》，第 IX 章）：按定义，这是指与此双线性型相伴的线性映射 \(\mathrm{X}(\mathrm{S}) \to \mathrm{Hom}_{\mathbf{Z}}(\varGamma(\mathrm{S}), \mathbf{Z})\) 与 \(\varGamma(\mathrm{S}) \to \mathrm{Hom}_{\mathbf{Z}}(\mathrm{X}(\mathrm{S}), \mathbf{Z})\) 都是双射。

显然，若命题的结论对两个环面成立，则对它们的积也成立。于是，由于每个 n 维环面同构于 \(U^{n}\) ，我们可化到 S = U 的情形。在这一特殊情形，断言是直接的。

设 \(f: S \to S'\) 是环面的态射。则线性映射 \(X(f): X(S') \to X(S)\) 与 \(\Gamma(f): \Gamma(S) \to \Gamma(S')\) 中的每一个都是另一个的转置：对一切 \(a' \in X(S')\) 与一切 \(X \in \Gamma(S)\) ，

\[
\langle X (f) (a ^ {\prime}), X \rangle = \langle a ^ {\prime}, \Gamma (f) (X) \rangle .\tag{3}
\]

命题 4. 设 S 与 S' 是环面。用 M(S,S') 表示从 S 到 S' 的李群态射所成的群。映射 \(f \mapsto \mathrm{X}(f)\) 与 \(f \mapsto \Gamma(f)\) 分别是从 M(S,S') 到 Hom \(_{\mathbf{Z}}(\mathrm{X}(\mathrm{S}'),\mathrm{X}(\mathrm{S}))\) 与到 Hom \(_{\mathbf{Z}}(\Gamma(\mathrm{S}),\Gamma(\mathrm{S}'))\) 的群同构。

若 \(f\) 是从 S 到 \(S'\) 的李群态射，则同态 \(X(f)\) 就是《谱理论》第 II 章 §1 第 7 节意义下 \(f\) 的对偶。同前所定义的映射 \(\varphi \mapsto \hat{\varphi}\) （从 \(\mathrm{Hom}_{\mathbf{Z}}(\mathrm{X}(\mathrm{S}'), \mathrm{X}(\mathrm{S}))\) 到 \(\mathrm{M}(\mathrm{S},\mathrm{S}')\) ）是映射 \(f \mapsto \mathrm{X}(f)\) （从 \(\mathrm{M}(\mathrm{S},\mathrm{S}')\) 到 \(\mathrm{Hom}_{\mathbf{Z}}(\mathrm{X}(\mathrm{S}'), \mathrm{X}(\mathrm{S}))\) ）的逆；因而后者是双射。若把 \(\Gamma(\mathrm{S})\) （相应地，\(\Gamma(\mathrm{S}')\) ）与 \(\mathbf{Z}\) -模 \(\mathrm{X}(\mathrm{S})\) （相应地，\(\mathrm{X}(\mathrm{S}')\) ）的对偶等同起来（命题 3），则 \(\Gamma(f)\) 与同态 \(X(f)\) 的转置相合，命题得证。

注。1) 设 \(f: \mathrm{S} \to \mathrm{S}'\) 是环面的态射。与 (1) 相伴的蛇形图（《代数学》，第 X 章，§1，第 2 节）给出正合列

\[
0 \longrightarrow \mathrm{Ker} \Gamma (f) \longrightarrow \mathrm{KerL} (f) \longrightarrow \mathrm{Ker} f \stackrel {d} {\longrightarrow}\tag{4}
\]

\[
\stackrel {{d}} {{\longrightarrow}} \operatorname{Coker} \Gamma (f) \longrightarrow \operatorname{Coker} \mathrm{L} (f) \longrightarrow \operatorname{Coker} f \longrightarrow 0.
\]

特别地，假定 \(f\) 满射且核为有限群 \(\mathbf{N}\) ，于是有正合列

\[
0 \longrightarrow \mathrm{N} \stackrel {i} {\longrightarrow} \mathrm{S} \stackrel {f} {\longrightarrow} \mathrm{S} ^ {\prime} \longrightarrow 0,
\]

其中 i 是典则单射。此时 \(L(f)\) 是双射，且 (4) 给出一个同构 \(N \to \text{Coker } \Gamma(f)\) ，从而有正合列

\[
0 \longrightarrow \Gamma (\mathrm{S}) \stackrel {\Gamma (f)} {\longrightarrow} \Gamma (\mathrm{S} ^ {\prime}) \longrightarrow \mathrm{N} \longrightarrow 0.\tag{5}
\]

此外，由《谱理论》第 II 章 §1 第 7 节定理 4，序列

\[
0 \longrightarrow \mathrm{X} (\mathrm{S} ^ {\prime}) \stackrel {{\mathrm{X} (f)}} {{\longrightarrow}} \mathrm{X} (\mathrm{S}) \stackrel {{\mathrm{X} (i)}} {{\longrightarrow}} \mathrm{X} (\mathrm{N}) \longrightarrow 0\tag{6}
\]

是正合的。

\begin{enumerate}
\def\labelenumi{\arabic{enumi})}
\setcounter{enumi}{1}
\tightlist
\item
  由命题 4，映射 \(f \mapsto \Gamma(f)(2\pi i)\) （从 \(\mathrm{M}(\mathbf{U},\mathrm{S})\) 到 \(\Gamma(\mathrm{S})\) ）是同构；若 \(a \in \mathrm{X}(\mathrm{S}) = \mathrm{M}(\mathrm{S},\mathbf{U})\) 且 \(f \in \mathrm{M}(\mathbf{U},\mathrm{S})\) ，则复合 \(a \circ f \in \mathrm{M}(\mathbf{U},\mathbf{U})\) 是自同态 \(u \mapsto u^r\) ，其中 \(r = \langle a,\Gamma(f)(2\pi i)\rangle\) 。我们从现在起把 \(\mathrm{M}(\mathbf{U},\mathbf{U}) = \mathrm{X}(\mathbf{U})\) 与 \(\mathbf{Z}\) 等同起来，其中元素 \(r\) （\(\mathbf{Z}\) 中的）对应自同态 \(u \mapsto u^r\) ；于是，采用上面的记号，
\end{enumerate}

\[
a \circ f = \langle a, \Gamma (f) (2 \pi i) \rangle .
\]

\begin{enumerate}
\def\labelenumi{\arabic{enumi})}
\setcounter{enumi}{2}
\item
  对正合列 \(0 \to \Gamma(\mathrm{S}) \to \mathrm{L}(\mathrm{S}) \xrightarrow{\exp_{\mathrm{S}}} \mathrm{S} \to 0\) 联系着一个从 \(\Gamma(\mathrm{S})\) 到 \(\mathrm{S}\) 的基本群上的同构，以下称之为典则的。对环面的任一态射 \(f: \mathrm{S} \to \mathrm{S}'\) ，\(\Gamma(f)\) 便可借助典则同构 \(\Gamma(\mathrm{S}) \to \pi_1(\mathrm{S})\) 与 \(\Gamma(\mathrm{S}') \to \pi_1(\mathrm{S}')\) 而与同态 \(\pi_1(f): \pi_1(\mathrm{S}) \to \pi_1(\mathrm{S}')\) （由 \(f\) 诱导的）等同起来。注意这给出了正合列 (5) 的另一种解释（参见《一般拓扑学》，第 XI 章，编写中）。
\item
  \(\mathbf{Z}\) -模同态 \(\delta : \mathrm{X}(\mathrm{S}) \to \mathrm{Hom}_{\mathbf{C}}(\mathrm{L}(\mathrm{S})_{(\mathbf{C})}, \mathbf{C})\) 与 \(\iota : \Gamma(\mathrm{S}) \to \mathrm{L}(\mathrm{S})_{(\mathbf{C})}\) （\(\iota\) 由 \(\Gamma(\mathrm{S})\) 到 \(\mathrm{L}(\mathrm{S})\) 中的典则单射诱导）扩张成 \(\mathbf{C}\) -向量空间的同构
\end{enumerate}

\[
\begin{array}{l} u: \mathbf {C} \otimes \mathrm{X} (\mathrm{S}) \to \mathrm{Hom} _ {\mathbf {C}} (\mathrm{L} (\mathrm{S}) _ {(\mathbf {C})}, \mathbf {C}) \\ v: \mathbf {C} \otimes \Gamma (\mathrm{S}) \to \mathrm{L} (\mathrm{S}) _ {(\mathbf {C})} \end{array}
\]

以下我们称之为典则的。注意，若把 \(X(S)\) 与 \(\Gamma(S)\) 之间的配对以 C-线性性扩张成双线性型 \(\ll\) , \(\gg\) （在 \((\mathbf{C} \otimes \mathrm{X}(S)) \times (\mathbf{C} \otimes \Gamma(S))\) 上），则

\[
\langle u (a), v (b) \rangle = 2 \pi i \ll a, b \gg .
\]

